\documentclass[10pt,a4paper]{amsart}
\usepackage[margin=19mm]{geometry}
\usepackage{amsmath,amssymb,mathtools}
\usepackage[scr=rsfs,cal=euler]{mathalfa}
\usepackage{xfrac}
\usepackage[unicode,hidelinks,bookmarks=false]{hyperref}
\newcommand{\ie}{i.e.\ }
\newcommand{\cf}{cf.\ }
\newcommand{\resp}{respectively\ }
\newcommand{\Subsection}{\subsection}
\providecommand{\nobreakdash}{\nobreak\hbox{-}}
\setlength{\emergencystretch}{2em}
\multlinegap0pt
\usepackage{amssymb}
\newtheorem{proposition}{Proposition}
\newcommand{\wind}[2]{[#1]_{#2}^{\sigma}}
\title{Iterated Hopf Ore Extensions over Group Rings}
\author{Can Hatipo\u{g}lu, Christian Lomp}
\date{Source: arXiv:2602.10850v2 (CC BY 4.0)}
\begin{document}
\maketitle

\noindent\emph{Source and licence.} Iterated Hopf Ore Extensions over Group Rings. Version arXiv:2602.10850v2 (CC BY 4.0), distributed by arXiv under the Creative Commons Attribution 4.0 International licence, http://creativecommons.org/licenses/by/4.0/, as recorded in the arXiv licence metadata for this exact version and shown on the abstract page https://arxiv.org/abs/2602.10850v2. Full licence notice: \texttt{licenses/arxiv-2602.10850-CC-BY-4.0.txt}.\par\smallskip

\noindent\textit{Editorial scope.} Counted unit: the article's proposition prop:iso-V-lambda-mu-rho (source0.tex lines 802--805). The article's complete original proof of it is delivered with it (source0.tex lines 806--874, one \texttt{proof} environment of 69 lines). No mathematics is written, repaired or re-derived: the statement and the proof are the article's own lines, in the article's own notation, and no retained line is substituted or re-flowed.  No further source-local context is needed: every cross-reference of every retained line resolves inside the two delivered spans.  Nothing was omitted from any retained span, so the package carries no omission record.  No retained line contains a citation, so the wrapper carries no bibliography and no bibliography entry.  No retained line contains a figure environment, an image-inclusion command or drawing code, so no image is delivered and the package declares no asset dependency.  Editorial formatting only, in addition: an AMS \texttt{amsart} preamble that restates the article's own declarations and macros used by the retained lines, namely the environments \texttt{proposition} on separate counters, together with the standard packages those lines need.  The retained environments print in this document as Proposition 1 in order of appearance, whereas the article prints the same items with the numbering of its own full document.  The complete native source of the article as distributed in the arXiv e-print for this exact version is bundled unchanged as \texttt{sources/194436/source0.tex}.
\begin{proposition}\label{prop:iso-V-lambda-mu-rho}
$V_x(\rho_1, \lambda_1,\mu_1) \simeq V_x (\rho_2, \lambda_2,\mu_2)$ as left $H$-modules if and only if there exists $0\leq k <n$ such that 
    $$\lambda_1=\lambda_2, \qquad \rho_1=\rho_2\chi^{-k},  \qquad \mu_1 = \mu_2 + \rho_2(\wind{e}{k}).$$
\end{proposition}

\begin{proof}
($\Rightarrow$) 
Set $V_1=V_x(\rho_1, \lambda_1,\mu_1)$ and $V_2=V_x (\rho_2, \lambda_2,\mu_2)$ and suppose  $\varphi:V_1\to V_2$ is an $H$-module isomorphism. Let $\{v_0,\ldots, v_{n-1}\}$ be a basis of $V_1$ and $\{w_0, \ldots, w_{n-1}\}$ be a basis of $V_2$. Consider 
%
\[
\varphi(v_0) = \sum_{i=0}^{n-1} k_i w_i.
\]
%
For any $g\in G$ we obtain
%
\[
\sum_{i=0}^{n-1} k_i \rho_1(g) w_i = \varphi(gv_0) = g\varphi(v_0) = \sum_{i=0}^{n-1} k_i gw_i 
= \sum_{i=0}^{n-1} k_i \rho_2(\sigma^{-i}(g)) w_i 
\]
%
Thus, for any $i$ with $k_i\neq 0$ we conclude $\rho_1(g) = \rho_2(\sigma^{-i}(g)) = \rho_2(g)\chi^{-i}(g)$. Since $\chi$ has order $n$ and since $\varphi$ is injective, there exists precisely one index $i$ such that $k_i\neq 0$ and we conclude $\rho_1 = \rho_2 \chi^{-i}$ as well as $\varphi(v_0)= k_i w_i$. Furthermore, 
%
\[ \lambda_1 k_i w_i = \varphi(x^n v_0) = x^n \varphi(v_0) =k_i \lambda_2 w_i.\]
%
Hence $\lambda_1=\lambda_2 =: \lambda $.  For $i>0$ we have 
%
\begin{eqnarray*}
 y\varphi(v_0) &=& k_i yw_i = k_i (\mu_2 + \rho(\wind{e}{i}))w_{i-1},\\
 \varphi(yv_0) &=& \lambda^{-1}\mu_1 \varphi(v_{n-1})  = \lambda^{-1}\mu_1 k_i x^{n-1} w_i   = k_i \mu_1 w_{i-1} .
 \end{eqnarray*}
%
Thus $\mu_1 = \mu_2 + \rho(\wind{e}{i})$. For $i=0$ we have 
%
\begin{eqnarray*}
 y\varphi(v_0) &=& k_0 y w_0 = k_0\lambda^{-1}\mu_2 w_{n-1} \\
 \varphi(yv_0) &=& \lambda^{-1}\mu_1 \varphi(v_{n-1})  = \lambda^{-1}\mu_1 k_0 w_{n-1}.
 \end{eqnarray*}
Hence $\mu_1 = \mu_2 = \mu_2 + \rho(\wind{e}{0})$.

($\Leftarrow$) Conversely, suppose $\lambda_1=\lambda_2$ and that there exists $0\leq k < n$ such that  $\rho_1=\rho_2\chi^{-k} $ and $\mu_1 = \mu_2 + \rho_2(\wind{e}{k})$ hold. 
Define a linear map $\varphi: V_1 \to V_2$ by $\varphi(v_{i}) = x^i w_k$, for $0\leq i<n$. Then
\[ \varphi(xv_i) = \varphi(v_{i+1}) = x^{i+1}w_k = x \varphi(v_i),\]  for $i<n-1$,  while  $\varphi(xv_{n-1})=\lambda \varphi(v_0)= \lambda w_k = x x^{n-1} w_k= x\varphi(v_{n-1}).$

For  $g \in G$, we have, using $\rho_1 = \rho_2 \chi^{-k} = \rho_2\circ \sigma^{-k}$:
%
\[\varphi(g v_{i}) = \rho_1 (\sigma^{-i}(g)) \varphi(v_{i}) = \rho_1 (\sigma^{-i}(g)) x^iw_k = \rho_2 (\sigma^{-i-k}(g)) x^iw_k\]

%
while
%
\[g \varphi(v_{i}) = g x^i w_k = \chi^{-i}(g) x^i \rho_2(\sigma^{-k}(g)) w_k  =\rho_2(\sigma^{-i-k}(g)) \varphi(v_i).\]
%
Thus $\varphi$ is $G$-linear. To check the $y$-linearity, note
%
\[
\varphi(yv_0)=\lambda^{-1}\mu_1\varphi(v_{n-1})
=\mu_1\lambda^{-1}x^{n-1}w_k 
=(\mu_2 + \rho_2(\wind{e}{k})w_{k-1} = y w_k = y\varphi(v_0).
\]
%


Finally, for $i \geq 1$, we have
%
\[\varphi(y v_{i}) = \varphi((\mu_1 + \rho_1 (\wind{e}{i}))v_{i-1}) = (\mu_1 + \rho_1 (\wind{e}{i})) x^{i-1}w_{k}\]
%
while
%
\[y \varphi(v_{i}) = yx^i w_k = (x^iy + x^{i-1}\wind{e}{i})w_k = (\mu_2 + \rho_2 (\wind{e}{i+k})) x^{i-1} w_{k}\]
%
Since $\mu_1 = \mu_2 + \rho_2(\wind{e}{k})$ and $\rho_1=\rho_2\circ \sigma^{-k}$, we obtain
\[ \mu_1 + \rho_1 (\wind{e}{i}) = \mu_2 + \rho_2(\wind{e}{k} + \sigma^{-k} (\wind{e}{i})) = \mu_2 + \rho_2 (\wind{e}{i+k}).\]
Thus $\varphi$ is a non-zero homomorphism of $H$-modules and hence an isomorphism as $V_1$ and $V_2$ are simple.
\end{proof}

\end{document}
